\documentclass[11pt,a4paper]{article}
\usepackage[margin=22mm]{geometry}
\usepackage[T1]{fontenc}
\usepackage{lmodern,amsmath,amssymb,array,booktabs}
\usepackage[colorlinks=true,linkcolor=blue,urlcolor=blue]{hyperref}
\hypersetup{pdftitle={The MiniJava Type System - CS3300 course revision},pdfauthor={CS3300 course staff; based on Jens Palsberg (2011)}}
\usepackage{fancyhdr}
\pagestyle{fancy}\fancyhf{}
\fancyhead[L]{\small MiniJava Type System}
\fancyhead[R]{\small CS3300 / September 2026}
\fancyfoot[C]{\thepage}
\setlength{\headheight}{14pt}
\setlength{\parindent}{0pt}
\setlength{\parskip}{5pt}
\newcommand{\kw}[1]{\text{\texttt{#1}}}
\newcommand{\mt}{\operatorname{methodtype}}
\newcommand{\mn}{\operatorname{methodname}}
\newcommand{\cn}{\operatorname{classname}}
\newcommand{\fld}{\operatorname{fields}}
\newcommand{\dist}{\operatorname{distinct}}
\newcommand{\noov}{\operatorname{noOverloading}}
\newcommand{\id}{\mathit{id}}
\newcommand{\ty}{\mathit{t}}
\newcommand{\ruleeq}[3]{\begin{equation}\frac{\begin{gathered}#2\end{gathered}}{\begin{gathered}#3\end{gathered}}\tag{#1}\end{equation}}
\newcommand{\axiom}[2]{\begin{equation}#2\tag{#1}\end{equation}}
\newcommand{\method}{\kw{public}\ t\ \id^M(t_1^F\ \id_1^F,\ldots,t_n^F\ \id_n^F)\ \{t_1\ \id_1;\ldots;t_r\ \id_r;\ s_1\ldots s_q\ \kw{return}\ e;\}}
\begin{document}
\begin{center}
{\LARGE\bfseries The MiniJava Type System}\\[7pt]
{\large CS3300 course revision --- 7 September 2026}\\[5pt]
Based on Jens Palsberg's \emph{The MiniJava Type System}, October 2011
\end{center}

This course revision incorporates the Assignment 2 clarifications into the rules
 themselves. Rule (16) permits covariant return types in overriding; rule (21)
 permits a return expression to have a subtype of the declared return type.
 Section 4 explicitly includes reflexivity for all MiniJava types. Rule (15)
 remains unchanged. Original section and rule numbers are retained; rule (16)
 is split into (16a) and (16b). Method types consistently omit parameter names.

\section{What is MiniJava?}
MiniJava is a small Java-like language with classes, single inheritance, methods,
integer arrays, and basic statements and expressions. Method overloading is not
allowed.

The statement \texttt{System.out.println(...);} prints integers, and
\texttt{e.length} applies to an expression of type \texttt{int[]}.

\section{Syntax}
The core syntax below follows the original reference. The
\href{https://fplaunchpad.org/cs3300_m26/grammar/minijava/}{course MiniJava grammar}
 specifies the complete syntax accepted in Assignment 2, including its additional
 operators and statement forms. Nonterminals are italic, keywords are in
 typewriter font, and lists written with dots may be empty. Superscripts and
 subscripts distinguish metavariables.

\begingroup\small
\begin{align*}
\textit{Goal}\quad g &::= mc\ d_1\ldots d_n\\
\textit{MainClass}\quad mc &::= \kw{class}\ \id\ \{\kw{public\ static\ void\ main}\ (\kw{String[]}\ \id^S)\{\\
 &\qquad t_1\ \id_1;\ldots;t_r\ \id_r;\ s_1\ldots s_q\ \}\}\!\\
\textit{TypeDeclaration}\quad d &::= \kw{class}\ \id\ \{t_1\ \id_1;\ldots;t_f\ \id_f;\ m_1\ldots m_k\}\\
 &\mid \kw{class}\ \id\ \kw{extends}\ \id^P\ \{t_1\ \id_1;\ldots;t_f\ \id_f;\ m_1\ldots m_k\}\\
\textit{MethodDeclaration}\quad m &::= \kw{public}\ t\ \id^M(t_1^F\ \id_1^F,\ldots,t_n^F\ \id_n^F)\ \{\\
 &\qquad t_1\ \id_1;\ldots;t_r\ \id_r;\ s_1\ldots s_q\ \kw{return}\ e;\ \}\\
\textit{Type}\quad t &::= \kw{int[]}\mid\kw{boolean}\mid\kw{int}\mid\id\\
\textit{Statement}\quad s &::= \{s_1\ldots s_q\}\mid\id=e;\mid\id[e_1]=e_2;\\
 &\mid\kw{if}\ (e)\ s_1\ \kw{else}\ s_2\mid\kw{while}\ (e)\ s\\
 &\mid\kw{System.out.println}(e);\\
\textit{Expression}\quad e &::= p_1\mathbin{\kw{\&\&}}p_2\mid p_1<p_2\mid p_1+p_2\mid p_1-p_2\mid p_1*p_2\\
 &\mid p_1[p_2]\mid p.\kw{length}\mid p.\id(e_1,\ldots,e_n)\mid p\\
\textit{PrimaryExpression}\quad p &::= c\mid\kw{true}\mid\kw{false}\mid\id\mid\kw{this}\\
 &\mid\kw{new\ int}[e]\mid\kw{new}\ \id()\mid {!e}\mid(e)\\
\textit{IntegerLiteral}\quad c &::= \langle\textsc{integer literal}\rangle\\
\textit{Identifier}\quad\id &::= \langle\textsc{identifier}\rangle
\end{align*}
\endgroup
In a type position, $\id$ must name a declared class: for example,
\texttt{Animal} in \texttt{Animal pet;} is a class type.

\newpage
\section{Notation for Rules}
An inference rule has the form
\[
\frac{\textit{hypothesis}_1\quad\textit{hypothesis}_2\quad\cdots\quad\textit{hypothesis}_n}
{\textit{conclusion}}.
\]
It states that the conclusion follows whenever every hypothesis holds. A rule
with no hypotheses is an axiom; its horizontal bar may be omitted. A derivation
is a tree whose leaves are axioms and whose internal nodes apply inference rules.
The root is the judgment established by the derivation.

\section{Subtyping}
We write $t_1\leq t_2$ when $t_1$ is a subtype of $t_2$. The relation is reflexive
on \emph{all} MiniJava types. On class names, it is the reflexive, transitive
closure of the declared \texttt{extends} relation.
\axiom{1}{t\leq t}
\ruleeq{2}{t_1\leq t_2\qquad t_2\leq t_3}{t_1\leq t_3}
\ruleeq{3}{\kw{class}\ C\ \kw{extends}\ D\ \ldots\ \text{is in the program}}{C\leq D}
In particular, rule (1) includes
\[
\kw{int}\leq\kw{int},\qquad
\kw{boolean}\leq\kw{boolean},\qquad
\kw{int[]}\leq\kw{int[]}.
\]
There are no subtype relationships between these three types, or between any of
these types and a class type. Thus the only subtype of each of \texttt{int},
\texttt{boolean}, and \texttt{int[]} is itself.

\section{Type Environments}
A type environment $A$ is a finite mapping from identifiers to types; its domain
is written $\operatorname{dom}(A)$. For pairwise distinct identifiers,
$[\id_1:t_1,\ldots,\id_r:t_r]$ maps $\id_i$ to $t_i$ for $i\in1..r$.
Environment composition gives the right-hand environment precedence:
\axiom{4}{(A_1\cdot A_2)(\id)=
\begin{cases}A_2(\id)&\text{if }\id\in\operatorname{dom}(A_2),\\
A_1(\id)&\text{otherwise.}\end{cases}}

\newpage
\section{Helper Functions}
\subsection{The \textit{classname} Helper Function}
The function $\cn$ returns the name of a class declaration.
\axiom{5}{\cn(\kw{class}\ \id\ \{\kw{public\ static\ void\ main}(\kw{String[]}\ \id^S)\{\ldots\}\})=\id}
\axiom{6}{\cn(\kw{class}\ \id\ \{t_1\ \id_1;\ldots;t_f\ \id_f;\ m_1\ldots m_k\})=\id}
\axiom{7}{\cn(\kw{class}\ \id\ \kw{extends}\ \id^P\ \{t_1\ \id_1;\ldots;t_f\ \id_f;\ m_1\ldots m_k\})=\id}
\subsection{The \textit{methodname} Helper Function}
The function $\mn$ returns the name of a method declaration.
\axiom{8}{\mn(\kw{public}\ t\ \id^M(\ldots)\ \ldots)=\id^M}
\subsection{The \textit{distinct} Helper Function}
The predicate $\dist$ checks that identifiers in a list are pairwise distinct.
\ruleeq{9}{\forall i\in1..n:\ \forall j\in1..n:\ \id_i=\id_j\Rightarrow i=j}
{\dist(\id_1,\ldots,\id_n)}
\subsection{The \textit{fields} Helper Function}
The environment $\fld(C)$ contains the fields of $C$ and its superclasses.
Fields declared in $C$ take precedence over inherited fields of the same name.
\ruleeq{10}{\kw{class}\ \id\ \{t_1\ \id_1;\ldots;t_f\ \id_f;\ m_1\ldots m_k\}\ \text{is in the program}}
{\fld(\id)=[\id_1:t_1,\ldots,\id_f:t_f]}
\ruleeq{11}{\kw{class}\ \id\ \kw{extends}\ \id^P\ \{t_1\ \id_1;\ldots;t_f\ \id_f;\ m_1\ldots m_k\}\ \text{is in the program}}
{\fld(\id)=\fld(\id^P)\cdot[\id_1:t_1,\ldots,\id_f:t_f]}

\newpage
\subsection{The \textit{methodtype} Helper Function}
The value $\mt(\id,\id^M)$ is the ordered list of parameter types and return
type of method $\id^M$ in class $\id$, including inherited methods. It is
$\bot$ if no such method exists. A method type has the form
$(t_1,\ldots,t_n)\to t$. Parameter names are not part of the method type.
\begingroup\small
\ruleeq{12}{\kw{class}\ \id\ \{\ldots\ m_1\ldots m_k\}\ \text{is in the program}\\
\text{for some }j\in1..k:\ \mn(m_j)=\id^M\\
m_j\text{ is of the form }\method}
{\mt(\id,\id^M)=(t_1^F,\ldots,t_n^F)\to t}
\ruleeq{13}{\kw{class}\ \id\ \{\ldots\ m_1\ldots m_k\}\ \text{is in the program}\\
\text{for all }j\in1..k:\ \mn(m_j)\neq\id^M}
{\mt(\id,\id^M)=\bot}
\ruleeq{14}{\kw{class}\ \id\ \kw{extends}\ \id^P\ \{\ldots\ m_1\ldots m_k\}\ \text{is in the program}\\
\text{for some }j\in1..k:\ \mn(m_j)=\id^M\\
m_j\text{ is of the form }\method}
{\mt(\id,\id^M)=(t_1^F,\ldots,t_n^F)\to t}
\ruleeq{15}{\kw{class}\ \id\ \kw{extends}\ \id^P\ \{\ldots\ m_1\ldots m_k\}\ \text{is in the program}\\
\text{for all }j\in1..k:\ \mn(m_j)\neq\id^M}
{\mt(\id,\id^M)=\mt(\id^P,\id^M)}
\endgroup
Rule (15) is an inheritance lookup rule: it applies only when the subclass does
not declare the method. Its equality is unchanged.

\subsection{The \textit{noOverloading} Helper Function}
For a method declared in class $\id$ with direct superclass $\id^P$, either
there is no inherited method of that name, or the declaration must be a valid
override. An override has exactly the same parameter types, in the same order,
and a return type that is the same as, or a subtype of, the inherited return type.
\ruleeq{16a}{\mt(\id^P,\id^M)=\bot}{\noov(\id,\id^P,\id^M)}
\ruleeq{16b}{\mt(\id,\id^M)=(t_1,\ldots,t_n)\to t\\
\mt(\id^P,\id^M)=(t'_1,\ldots,t'_n)\to t'\\
t_i=t'_i\quad(i\in1..n)\qquad t\leq t'}{\noov(\id,\id^P,\id^M)}
Both parameter lists have length $n$; parameter names may differ. Changing a
parameter type, even to a subtype, or changing the parameter count is forbidden.

\newpage
\section{Type Rules}
\subsection{Type Judgments}
We use the following seven forms of judgments:
\begin{center}
\begin{tabular}{ll}
\toprule
Judgment & Meaning\\\midrule
$\vdash g$ & The goal (whole program) type checks.\\
$\vdash mc$ & The main class type checks.\\
$\vdash d$ & The type declaration type checks.\\
$C\vdash m$ & The method declaration type checks in class $C$.\\
$A,C\vdash s$ & The statement type checks in environment $A$, class $C$.\\
$A,C\vdash e:t$ & The expression has type $t$ in environment $A$, class $C$.\\
$A,C\vdash p:t$ & The primary expression has type $t$ in that context.\\
\bottomrule
\end{tabular}
\end{center}
\subsection{Goal}
\ruleeq{17}{\dist(\cn(mc),\cn(d_1),\ldots,\cn(d_n))\\
\vdash mc\qquad\vdash d_i\quad(i\in1..n)}{\vdash mc\ d_1\ldots d_n}
\subsection{Main Class}
\begingroup\small
\ruleeq{18}{\dist(\id_1,\ldots,\id_r)\qquad
[\id_1:t_1,\ldots,\id_r:t_r],\bot\vdash s_i\quad(i\in1..q)}
{\vdash\kw{class}\ \id\ \{\kw{public\ static\ void\ main}(\kw{String[]}\ \id^S)\{t_1\ \id_1;\ldots;t_r\ \id_r;\ s_1\ldots s_q\}\}}
\endgroup
Here $\bot$ in the class position denotes the absence of an instance context;
\texttt{this} is unavailable in the static main method.

\subsubsection*{Examples of the corrected return and overriding rules}
Assume \texttt{B extends A}. Each row considers the stated type constraint;
the rest of the program must also type check.
\begin{center}\small
\begin{tabular}{@{}p{0.71\linewidth}p{0.23\linewidth}@{}}
\toprule
Situation & Result\\\midrule
A method declared to return \texttt{A} returns \texttt{new B()}. & Allowed\\[5pt]
A method declared to return \texttt{B} returns \texttt{new A()}. & Type error\\[5pt]
An inherited \texttt{A f(A x)} is overridden by \texttt{B f(A y)}. & Allowed\\[5pt]
An inherited \texttt{A f(A x)} is replaced by \texttt{B f(B x)}. & Type error: parameter type changed\\[5pt]
An inherited \texttt{B f(A x)} is overridden by \texttt{A f(A x)}. & Type error: return type widened\\
\bottomrule
\end{tabular}\end{center}
The parameter equality in rule (16b) concerns method \emph{declarations}.
At a method \emph{call}, rule (35) permits argument expressions whose types
are subtypes of the declared parameter types.

\newpage
\subsection{Type Declarations}
\ruleeq{19}{\dist(\id_1,\ldots,\id_f)\\
\dist(\mn(m_1),\ldots,\mn(m_k))\\
\id\vdash m_i\quad(i\in1..k)}
{\vdash\kw{class}\ \id\ \{t_1\ \id_1;\ldots;t_f\ \id_f;\ m_1\ldots m_k\}}
\ruleeq{20}{\dist(\id_1,\ldots,\id_f)\\
\dist(\mn(m_1),\ldots,\mn(m_k))\\
\noov(\id,\id^P,\mn(m_i))\qquad\id\vdash m_i\quad(i\in1..k)}
{\vdash\kw{class}\ \id\ \kw{extends}\ \id^P\ \{t_1\ \id_1;\ldots;t_f\ \id_f;\ m_1\ldots m_k\}}
\subsection{Method Declarations}
\begingroup\small
\ruleeq{21}{\dist(\id_1^F,\ldots,\id_n^F)\qquad\dist(\id_1,\ldots,\id_r)\\
A=\fld(C)\cdot[\id_1^F:t_1^F,\ldots,\id_n^F:t_n^F]\cdot[\id_1:t_1,\ldots,\id_r:t_r]\\
A,C\vdash s_i\quad(i\in1..q)\qquad A,C\vdash e:t'\qquad t'\leq t}
{C\vdash\method}
\endgroup
The return expression's type $t'$ may be a subtype of the declared return type
$t$. This check applies to every method, whether or not it overrides another.
\subsection{Statements}
\ruleeq{22}{A,C\vdash s_i\quad(i\in1..q)}{A,C\vdash\{s_1\ldots s_q\}}
\ruleeq{23}{A(\id)=t_1\qquad A,C\vdash e:t_2\qquad t_2\leq t_1}{A,C\vdash\id=e;}
\ruleeq{24}{A(\id)=\kw{int[]}\qquad A,C\vdash e_1:\kw{int}\qquad A,C\vdash e_2:\kw{int}}{A,C\vdash\id[e_1]=e_2;}
\ruleeq{25}{A,C\vdash e:\kw{boolean}\qquad A,C\vdash s_1\qquad A,C\vdash s_2}{A,C\vdash\kw{if}\ (e)\ s_1\ \kw{else}\ s_2}
\ruleeq{26}{A,C\vdash e:\kw{boolean}\qquad A,C\vdash s}{A,C\vdash\kw{while}\ (e)\ s}
\ruleeq{27}{A,C\vdash e:\kw{int}}{A,C\vdash\kw{System.out.println}(e);}

\newpage
\subsection{Expressions and Primary Expressions}
\begingroup
\setlength{\abovedisplayskip}{7pt}\setlength{\belowdisplayskip}{7pt}
\ruleeq{28}{A,C\vdash p_1:\kw{boolean}\qquad A,C\vdash p_2:\kw{boolean}}{A,C\vdash p_1\mathbin{\kw{\&\&}}p_2:\kw{boolean}}
\ruleeq{29}{A,C\vdash p_1:\kw{int}\qquad A,C\vdash p_2:\kw{int}}{A,C\vdash p_1<p_2:\kw{boolean}}
\ruleeq{30}{A,C\vdash p_1:\kw{int}\qquad A,C\vdash p_2:\kw{int}}{A,C\vdash p_1+p_2:\kw{int}}
\ruleeq{31}{A,C\vdash p_1:\kw{int}\qquad A,C\vdash p_2:\kw{int}}{A,C\vdash p_1-p_2:\kw{int}}
\ruleeq{32}{A,C\vdash p_1:\kw{int}\qquad A,C\vdash p_2:\kw{int}}{A,C\vdash p_1*p_2:\kw{int}}
\ruleeq{33}{A,C\vdash p_1:\kw{int[]}\qquad A,C\vdash p_2:\kw{int}}{A,C\vdash p_1[p_2]:\kw{int}}
\ruleeq{34}{A,C\vdash p:\kw{int[]}}{A,C\vdash p.\kw{length}:\kw{int}}
\ruleeq{35}{A,C\vdash p:D\qquad\mt(D,\id)=(t'_1,\ldots,t'_n)\to t\\
A,C\vdash e_i:t_i\qquad t_i\leq t'_i\quad(i\in1..n)}{A,C\vdash p.\id(e_1,\ldots,e_n):t}
\axiom{36}{A,C\vdash c:\kw{int}}
\axiom{37}{A,C\vdash\kw{true}:\kw{boolean}}
\axiom{38}{A,C\vdash\kw{false}:\kw{boolean}}
\ruleeq{39}{\id\in\operatorname{dom}(A)}{A,C\vdash\id:A(\id)}
\ruleeq{40}{C\neq\bot}{A,C\vdash\kw{this}:C}
\ruleeq{41}{A,C\vdash e:\kw{int}}{A,C\vdash\kw{new\ int}[e]:\kw{int[]}}
\axiom{42}{A,C\vdash\kw{new}\ \id():\id}
\ruleeq{43}{A,C\vdash e:\kw{boolean}}{A,C\vdash {!e}:\kw{boolean}}
\ruleeq{44}{A,C\vdash e:t}{A,C\vdash(e):t}
\endgroup
\end{document}
